\section{Types of Journalism in the Press Corps}

\subsection{Beat articles}

“Beat” refers to going deep into a journalist's area of expertise or interest. Hence, a beat article consists of a written piece directed towards a specific topic of \emph{relevance}.\footnote{Kovach, B., \& Rosenstiel, T. (2014). The Elements of Journalism: What Newspeople Should Know and the Public Should Expect (3rd ed.). Three Rivers Press.} Under the MUN press committee, the topic or “beat” will depend on the committee you are assigned to report on. To fully comprehend the beat of the committee, it is important to understand the active interests of countries during the debates and the opposing viewpoints of the blocs. To do this, you must be actively present in the committee you’re assigned in. The best beat storytelling comes from paying attention to details that are representative of higher power dynamics at play and putting into context the topics being discussed.\footnote{Scanlan, C. (2000). Reporting and Writing: Basics for the 21st Century. Oxford University Press.} In addition, to conduct beat journalism one must be able to grasp which stories deserve being told and which can be discarded. This is shown in the examples attached below.

In the committee, it is recommended to use the first half of the day to collect information by actively listening to the debates and use the second half to write the article. It is necessary that every reporter submits one work (could be an article or any other media type) at the end of each day of debate, besides other daily submissions such as interviews, press conferences and breaking news articles. It is encouraged that your beat article includes information from secondary sources to enhance credibility but of course must include proper citation.\footnote{Mencher, M. (2010). Melvin Mencher's News Reporting and Writing (12th ed.). McGraw-Hill.}

Lastly, it is important that your article demonstrates the viewpoint of the editorial line of your assigned news agency. If, for example, your magazine is of western bias, your works must carry the same biased worldview. Ultimately, their reporting methods, linguistic style, and visual aesthetic must be emulated in your works.

\subsection{Examples of Beat articles}

\begin{enumerate}
  \item Whitlock, C. (2019, December 9). At war with the truth. \emph{The Washington Post}. Craig Whitlock goes into detail of the ways in which the US government has lied about the war in Afghanistan.\footnote{Whitlock, C. (2019, December 9). At war with the truth. The Washington Post. \url{https://www.washingtonpost.com/graphics/2019/investigations/afghanistan-papers/afghanistan-war-co} nfidential-documents/}
  \item A global revolution in attitudes towards cannabis is under way. (2019, August 29). \emph{The Economist}. The Economist delves into how the public opinion is changing towards cannabis.\footnote{A global revolution in attitudes towards cannabis is under way. (2019, August 29). The Economist. \url{https://www.economist.com/international/2019/08/29/a-global-revolution-in-attitudes-towards-cannabis-is-under-way}}
\end{enumerate}

\subsection{Interview}

The second media type a press delegate can use is the interview. The interview requires in person authorization by chairs and your assigned committee's chairs. The ways to conduct the interview are at the delegate’s discretion, however the content of the interview must be of relevance to the conference and loyal to the magazine the delegate represents. An interview can be used for multiple purposes:

\begin{enumerate}
  \item Inform the public: inform about committee progress or delegate positions through the collection of basic factual information
  \item Shaping public opinion: Influence the views of delegates through, for example, direct quotes from other influential delegates
  \item Gain insider information: reveal conflicts or alliances behind negotiations, discover behind the scenes operations, etc.
  \item Entertainment: bringing personality or humour to the news coverage to engage more people
\end{enumerate}

To conduct an interview one must identify the delegate and request an interview, under chairs' authorization. Begin the interview with a previously defined set of questions- but don’t be scared to veer off script if unprecedented topics of interest come up. Before conducting the interview, the press delegate should be well informed about the topic in order to handle these possible impromptu questions. And finally, keep the interviews short to ensure that only necessary information is collected and provided.

\subsection{Examples of interviews}

\begin{enumerate}
  \item Uncommon Knowledge with Peter Robinson. (2022, June 6). \emph{Is the end of history over? — with Francis Fukuyama} [Video] Interview with Francis Fukuyama.\footnote{Uncommon Knowledge with Peter Robinson. (2022, June 6). Is the end of history over? — with Francis Fukuyama [Video]. YouTube. \url{https://www.youtube.com/watch?v=czcpT-OeQJ4}}
  \item DW News. (2020, March 6). \emph{Interview with economist Thomas Piketty on capital and ideology} [Video] Interview with economist Thomas Piketty on capital and ideology.\footnote{DW News. (2020, March 6). Interview with economist Thomas Piketty on capital and ideology [Video]. YouTube. \url{https://www.youtube.com/watch?v=utnpOqE6Lk0}}
  \item BBC News. (2024). \emph{Joe Biden interview} [Video] Interview with Joe Biden conducted by the BBC.\footnote{BBC News. (2024). Joe Biden interview [Video]. YouTube. \url{https://www.youtube.com/watch?v=INuwqHzBFC0}}
\end{enumerate}

\subsection{Opinion Editorials (aka Op-Ed)}

In contrast to a beat article, in Op-Eds the reveal of the journalist’s and news agency’s point of view is expected.\footnote{Shipley, D. (2004, September 3). And now a word from Op-Ed. The New York Times. \url{https://www.nytimes.com/2004/02/01/opinion/and-now-a-word-from-op-ed.html}} These opinionated articles can be based on the agenda of the committee discussion. Op-Ed’s are powerful and interesting as they allow delegates to express subjective perspectives on events and discussions that occur in committees. However, they may be challenging because a well-formed opinion requires thorough research to defend it.\footnote{The Duke University Op-Ed Project. (n.d.). How to write an Op-Ed or column. \url{https://styleguide.duke.edu/toolkits/writing-media/how-to-write-an-op-ed-or-column/}} For that reason, delegates are expected to conduct this type of investigation for their Op-Ed works. The purpose of an Op-Ed can be to influence delegate perception, spark debate, showcase bias or alignment of a specific press outlet, etc.

Possible types of Op-Eds:

\begin{enumerate}
  \item Stance critiques: critique a certain delegate’s stance on a given matter through an alternative proposition or a rebuttal to what was said in committee discussion
  \item Historical context: providing a background or historical perspective to a contemporary issue of relevance
  \item Bias piece: A biased piece that reflects the bias of the assigned magazine or a piece that highlights possible bias of the committee discussion by highlighting omitted information
  \item Emotional Appeal: taking emotional angles on issues in order to rile up public support
  \item Satirical Op-Eds: using humour for social criticism
  \item Alternative perspective: a piece that suggests alternative solutions or perspectives to what was offered in the committee
\end{enumerate}

Things to consider while writing an Op-Ed:

\begin{itemize}
  \item Less is more: the simpler you can communicate your thoughts, the better. The message comes across in a more simple manner- making it more compelling.
  \item Op-Eds express opinions that are backed up by evidence, not baseless arguments
  \item Structure is key: all elements of Op-Ed must lead to a conclusion rather than leaving things open ended
\end{itemize}

\subsection{Examples of Op-Eds}

\begin{enumerate}
  \item Tagliapietra, S., \& Trasi, C. (2024, November). Trump's comeback and its implications for EU climate and energy policy. Bruegel. Simone Tagliapietra and Cecilia Trasi argue that “Trump’s return should not be seen as a threat to EU decarbonisation, but as a rallying cry to unite and push forward more strongly.”\footnote{Tagliapietra, S., \& Trasi, C. (2024, November). Trump's comeback and its implications for EU climate and energy policy. Bruegel. \url{https://www.bruegel.org/first-glance/trumps-comeback-and-its-implications-eu-climate-and-energy-poli} cy}
  \item PS Editors. (2024, November). Trump's return: PS commentators explain what it means. \emph{Project Syndicate}. A roundtable article that displays the opinion of multiple authors on what Trump's return to office can imply for the future of international politics.\footnote{PS Editors. (2024, November). Trump's return: PS commentators explain what it means. Project Syndicate. \url{https://www.project-syndicate.org/onpoint/trump-return-ps-commentators-explain-what-it-means-by-ps-editors-2024-11}}
\end{enumerate}

\subsection{Feature}

Features are arguably the most creative media type as they can adopt different storytelling structures: poetry, prose, letters, diary entries, monologue/internal thoughts (ex. \emph{Inside the mind of the delegate of China),} script/play format, fable, etc.\footnote{Hart, J. (2011). Storycraft: The Complete Guide to Writing Narrative Nonfiction. University of Chicago Press.} In a feature, storytelling takes precedence over research, however, the chosen topic must be relevant to the committee and research is still necessary in order to execute it well.\footnote{Ricketson, M. (2004). Writing Feature Stories: How to Research and Write Articles — From Listicles to Longform. Allen \& Unwin.} Features are best when they are really creative, so let your imagination take over.

\subsection{Examples of Features}

\begin{enumerate}
  \item \url{https://www.theguardian.com/us-news/ng-interactive/2025/jun/08/boy-geniuses-great-men-trump} \href{https://www.theguardian.com/us-news/ng-interactive/2025/jun/08/boy-geniuses-great-men-trump}{uses-great-men-trump} Argumentative essay on how the narrative for “smart men” perpetuates an obscured gender inequality.
  \item \url{https://poets.org/poem/september-1-1939} Widely-known poem written by
\end{enumerate}

W.H. Auden during the outbreak of World War ll. The poem captures feelings of fear and uncertainty in the face of fascism and war—as well as glimmers of hope that people might come together to counter authoritarians.

Now we are departing the more traditional forms of journalism and entering a newer frontier for reporting on developments - using audiovisual methods. With the rise of the Web, technological expansions of data speeds, and the invention of social media, more dynamic forms of reporting have now largely replaced traditional media for the attention spans demanded of the general public.\footnote{Newman, N. (2025). Journalism, Media, and Technology Trends and Predictions 2025. Reuters Institute for the Study of Journalism, University of Oxford. \url{https://reutersinstitute.politics.ox.ac.uk/journalism-media-and-technology-trends-and-predictions-2025}} Whilst there are significant overlaps in AV reporting and other forms of reporting, we highly encourage reporters to give these versions a try if they wish for their journals to look more of the times, so-to-speak.

For any AV submission, as we wish to publish this to the IEUMUN Press Corp’s Instagram account, they must be submitted in a publishable format. Therefore, videos must be sent in the size/ratio required for a story/reel (i.e., 9:16 or 1080x1920 in vertical format) and audio submissions must be sent in a video format, with graphics attached to the raw audio.

\subsection{Video Bulletin}

The first form of this is in short-form video bulletins. Bulletins are typically collations of smaller reports made by a journalist or an organisation. They are a great example of which being the evening news reports seen on television.\footnote{Boyd, A., Stewart, P., \& Alexander, R. (2012). Broadcast Journalism: Techniques of Radio and Television News (7th ed.). Routledge.} With a written newspaper, much of the description of the events come from describing the commotion that goes on with flowery language and other literary devices. With video reports in a bulletin, the viewer does not have to imagine but can see what goes on for themselves, and come to their own conclusions. With video, like with the written form, the journalist can choose what to show and in what light.\footnote{Zettl, H. (2016). Television Production Handbook (12th ed.). Cengage Learning.} If you choose to go down this route, record either as much footage as you can to edit down, or pick and choose speakers or shots that help perpetuate your journal’s leanings or biases.

When we look for a video submission, we expect either a single debate covered in depth or a wide range of topics covered, even if a round-up with some form of interpretation of each of them. For a video bulletin, we will expect the latter. Whilst we won’t put a hard-and-fast rule on the minimum number of topics discussed in each, we expect journalists to not stick to too similar topics or only discuss the happenings of only a few at a surface level.

\emph{Examples of Video Bulletins:}

\begin{enumerate}
  \item BBC News. (2021, December 7). \emph{BBC News at Ten} [Video] The \emph{BBC News}’ evening bulletin in December 2021, discussing an at-the-time alleged party during the COVID-19 pandemic, the British Armed Forces’ withdrawal from Afghanistan, and coverage of cricket.\footnote{BBC News. (2021, December 7). BBC News at Ten [Video]. YouTube. \url{https://www.youtube.com/watch?v=pAOxtrqZSnA}}
  \item \emph{Euronews. (2025, June 3). Latest news bulletin | June 3rd – Morning [Video]. Euronews}’ roundup and look-forward in June 2025, covering failed peace talks between Russia and Ukraine held in Turkey, and the eruption of Mount Etna, the largest active volcano in Europe.\footnote{Euronews. (2025, June 3). Latest news bulletin | June 3rd – Morning [Video]. YouTube. \url{https://www.youtube.com/watch?v=pjMolFT_EyA}}
  \item BBC Three. (2016, January 4). \emph{60 Seconds} [Video] An example of a video bulletin focussed on short form and a target audience in mind (BBC Three’s traditional audience are ages 16-35.) Example discusses what international broadcast news had been reporting in January 2016, including recent developments of the Islamic State and large-scale flooding in southern India.\footnote{BBC Three. (2016, January 4). 60 Seconds [Video]. YouTube. \url{https://www.youtube.com/watch?v=NekVSWCTTac}}
\end{enumerate}

\subsection{Podcast}

Podcasts are one of the more popular forms of journalism in the current day. Podcasting is ostensibly the distribution of audio programs for download over the internet, distributing traditionally over RSS feeds or on video-sharing platforms, such as YouTube.\footnote{Berry, R. (2016). Podcasting: Considering the evolution of the medium and its association with the word "radio." The Radio Journal: International Studies in Broadcast \& Audio Media, 14(1), 7–22. \url{https://doi.org/10.1386/rjao.14.1.7_1}} Originally conceived by Apple to be incorporated into their iTunes platform for iPods (i\textbf{Pod} broad\textbf{casts}), and ostensibly a more modern take on radio and television talk shows, the accessibility and relative ease the format lends hosts has meant the medium has skyrocketed in popularity with estimates suggesting over 200 million episodes of podcasts have been published by June 2024.\footnote{Podcast Index. (2024). Podcast statistics. \url{https://podcastindex.org/}; Newman, N. (2025). Journalism, Media, and Technology Trends and Predictions 2025. Reuters Institute.} Additionally, due to the majority of podcasts being accessible for free and with many providers being available that provide podcast streams, popularity of the medium has never been as high.

Formats of podcasts are, by its own nature, quite open to interpretation. Episodes can follow rigid, scripted structures where it, in effect, acts as audio editions or compendiums to other written works. Other times, they can have very little structure and often act as a vehicle for free flowing conversation, discussion, or debate between participants.\footnote{Lindgren, M. (2016). Personal narrative journalism and podcasting. The Radio Journal: International Studies in Broadcast \& Audio Media, 14(1), 23–41.} However, almost all podcasts have at least one central figure (whether it be an individual or a group) that episodes use as a host, or at the very least as a guiding figure to move the episode along. A good amount of podcasts will feature guests or other contributors to add additional context to stories or provide contrasting opinions. In addition, many podcasts have supplementary content for those who wish to learn more (such as show notes, references, biographies, or transcripts.) Within the Press Corps, supplementary info is not expected for any podcast to be made.

As journalists, podcasts can be an exciting avenue to pursue as you can discuss, explain, examine, invite whoever, whatever you can, and it can allow you to portray your side of any argument in the depth the structure you choose can provide. Do beware, however, that an episode doesn’t fall into the trap of being unfocussed. When the format is so free, it can often feel like it has no structure or point. If you choose to produce a podcast, we highly recommend having some form of goal or angle you wish to get across, and build everything else around that.

\emph{Examples of Podcasts:}

\begin{enumerate}
  \item \emph{The Daily T Podcast}. (n.d.). \emph{The Daily Telegraph}. The Daily T Podcast is a compendium to the right-leaning British newspaper \emph{The Daily Telegraph}, and acts as a daily roundup of its headlines. It also serves as a good example of how to transform a traditional journalistic source into a more modern format and across from the written word to the spoken word.\footnote{The Daily T Podcast. (n.d.). The Daily Telegraph. \url{https://www.telegraph.co.uk/the-daily-t-podcast/}}
  \item \emph{The Morning Brief}. (n.d.). \emph{The Economist} [Video podcast] The Morning Brief is published by \emph{The Economist} as additional analysis of the world of business and politics, as well as forecasts to how financial markets will perform and its impacts on businesses and the general public.\footnote{The Morning Brief. (n.d.). The Economist [Video podcast]. YouTube. \url{https://www.youtube.com/playlist?list=PLmVg8_8GDpP7_BjENhqQdF5On0FIKIwKq}}
\end{enumerate}

\subsection{Documentary}

A documentary film is a long-form piece of work that centers in and around a central figure or idea and goes into quite a bit of depth, with the aim of providing some form of narrative or message for the end.\footnote{Nichols, B. (2017). Introduction to Documentary (3rd ed.). Indiana University Press.} Documentaries tend to follow similar act structures to scripted fictional films, with the aim to instruct and educate, or provide a record of the past for future watchers. Within the past, when the medium used to be known as an ‘actuality’ film, they were traditionally only a few minutes in length and showed as true a perspective of reality as possible.\footnote{Barnouw, E. (1993). Documentary: A History of the Non-Fiction Film (2nd rev. ed.). Oxford University Press.} However, actuality pictures often bored audiences and were seen as quite dry. Nowadays, runtimes have gotten longer and the boundaries for which a documentary can fit in have only grown.

Even though documentaries have existed for nearly as long as films have been around, more modern interpretations are pervasive in the medium that bend the more   traditional   rules   of   documentaries,   and   are   often one-person/very-small-team productions. A major example is the video essay.\footnote{Borum Chattoo, C., \& Feldman, L. (2020). A Comedian and an Activist Walk into a Bar: The Serious Role of Comedy in Social Justice. University of California Press} Popular on social media sites such as YouTube and Dailymotion, they act as deep dives on current or historical topics, themes, or issues and often follow similar narrative devices to traditional documentaries. However, video essayists use their new medium to its advantage with clever uses of video editing, audio, and more current cultural references to their advantage.

As the conference is only a few days in length, we don’t expect any budding documentarians to write, film, edit, and submit feature-length films. However, if journalists wish to make a longer piece (over five-or-so minutes), then this can be worked on over multiple sessions with smaller other pieces of journalism, sneak peeks, or status updates submitted at the end of each session in lieu. However, this must be submitted by the end of the penultimate session as the final collaborative journal will be in written form.

\emph{Examples of Documentaries:}

\begin{enumerate}
  \item Blue Planet II - A sequel to the 2001 show of the same name, the eight-volume series follows David Attenborough and his team to examine the ongoing state of Earth’s oceans and seas, alongside the continuing environmental damage that is being done to it. The series took over five years to make, with over 6,000 hours of footage recorded from over 4,000 dives.\footnote{Attenborough, D. (Presenter). (2017). Blue Planet II [Documentary series]. BBC Studios. \url{https://en.wikipedia.org/wiki/Blue_Planet_II}}
  \item A Canadian Tragedy: The Nortel Saga - A two-part video essay that deep-dives one of the most influential telecommunications companies of all time, its links to Canadian political history, and its downfall - with a retrospective look as to how its lessons have still to be learned in modern business practices. A fantastic example of how to use the YouTube format to your advantage, in terms of presentation.\footnote{A Canadian Tragedy: The Nortel Saga [Video essay]. YouTube. \url{https://www.youtube.com/playlist?list=PLAB-wWbHL7VsRxU4zhwkdsaiEh7NcL7XG}}
  \item Educating Yorkshire - A documentary series made in the early 2010s, it is an example of the documentary technique ‘fly on the wall’. This method involves minimal involvement from the production crew, watching their subjects go by their day as normal. This normally involves inconspicuous cameras and frequent interviews to discuss their actions. The series given is done to great effect showing the current state of British secondary education in northern England, one of the more deprived areas of the country.\footnote{Educating Yorkshire [Documentary series]. (2013). Channel 4. \url{https://www.youtube.com/playlist?list=PLGkph1NtiNFXuIZZL4eM4G6AtbSCI8aLL}}
\end{enumerate}

\subsection{Editorial Cartoons}

Fancy yourself an artist? Editorial cartoons, sometimes known as political cartoons, act as caricatures of current or historical affairs in order to satirise, promote, or demonise an argument - traditionally the artist’s own opinion.\footnote{Hess, S., \& Northrop, S. (2011). American Political Cartoons: The Evolution of a National Identity, 1754–2010. Transaction Publishers.} The origins of editorial cartoons go back to the invention of magazines and the newspaper, and often served many different points. Firstly was advertising for their publications; second was propaganda for either a political stance; lastly, it was seen as an accessible form of media absorption for people who weren’t perfectly literate.\footnote{Press, C. (1981). The Political Cartoon. Fairleigh Dickinson University Press.} Soon after, magazines such as \emph{Punch}, \emph{Puck}, and \emph{Mad} - purely dedicated to these comics - entered publication and quickly entered the zeitgeist, acting as antidotes to written journals and newspapers, which were seen as dry and oftentimes politically conservative.

In the modern era, political cartoons and caricatures tend to be found in two main areas. First is still in the printed form - magazines such as \emph{Mad} still live on strong and have often inspired newer ones to pop up such as \emph{Private Eye}. The other location where these can be discovered is online. With the new ease of creating digital art and publicising work (on platforms such as Twitter), artists of all political persuasions have found their home and audiences in the online realm, especially for political beliefs that are not as mainstream.\footnote{Walker, R. (2003). The Cartoon History of Time. HarperCollins; Lamb, C. (2004). Drawn to Extremes: The Use and Abuse of Editorial Cartoons. Columbia University Press.}

Size of art can range dramatically as well, with large intricate pieces that have taken weeks alongside ‘pocket cartoons’ - smaller, line-art style that can be made quickly and tend to be more topical to recent developments. Within the committee, if you would like to make a cartoon, we recommend keeping drawings to as small as you reasonably can make it. We don’t wish for journalists to spend multiple days on a single piece as this may not truly reflect a country’s position in committee come the time of publication. Pocket cartoons are our recommendation, but if you think you can do something more detailed in the span of a session, then be our guest!

A final small note, in keeping with the conference’s -- and our own -- expectations of generative AI usage (ChatGPT, Gemini, etc.), we will not accept any editorial cartoons that are created using these tools. Believe us, we can tell!

\emph{Examples of Editorial Cartoons:}

\guidefig{assets/figure-01.png}[0.36]

\guidefig[A \emph{New York Times} cartoon from 1903, portraying the United States attempting to influence Latin American politics, especially their impact on the separation of Colombia and Panama.\protect\footnotemark]{assets/figure-02.jpg}[0.29]\footnotetext{Crane, F. (1903, November 15). The man behind the egg. [Political cartoon]. Library of Congress Manuscript Division, Theodore Roosevelt Digital Library. Dickinson State University. \url{https://www.theodorerooseveltcenter.org/digital-library/o302256}}

Pocket cartoon by Matt Pritchett for \emph{The Daily Telegraph} in August 2023, satirising the reaction by Londoners to the decision to implement an ultralow emissions zone across almost all of the British capital.\footnote{Pritchett, M. (2023, August 31). It was owned by a little old lady who just used it to pop out and vandalise Ulez cameras. [Political cartoon]. Telegraph Newsprints. \url{https://telegraph.newsprints.co.uk/39851798-matt-cartoon-it-was-owned-by-a-little-old-lady-who-just-u} sed-it-to-pop-out-and-vandalise-ulez-cameras-347425880\_matt-pritchett\_matt-cartoon-31-august-202 3-art/}
